\section{Problem Statement}\label{sec:problem-statement}

Security operations are bottlenecked by the manual translation of operator intent
into tool actions. A penetration tester or analyst must map free-form
instructions---``scan the internal network for EternalBlue'', ``check this host
for CVE-2021-44228''---onto commands, targets and constraints while respecting
authorisation boundaries. Recent work shows that Large Language Models (LLMs) can
act as capable copilots for parts of this work, producing realistic attack plans
and driving tools in controlled environments \parencite{deng2024pentestgpt,
happe2023pwned}. Using an LLM for this translation step is attractive precisely
because natural language is the operator's native interface.

The problem is that a language model is not a reliable translator. Its output can
be malformed, misclassified, or wholly fabricated: LLMs hallucinate identifiers
and parameters \parencite{ji2023hallucination}, and they can be diverted by
adversarially crafted utterances \parencite{greshake2023indirect,
liu2023promptinjection}. In a purely conversational copilot these errors are
display text that a human can vet. In an \emph{automated} security pipeline,
where intent is converted directly into actions, each such error becomes a
candidate machine-dispatched action: a fabricated CVE is looked up, a wrong
target is scanned, or an injected instruction is obeyed. Acting on unvalidated
model output therefore transfers every residual error of the model into the
security toolchain it drives.

The intent-recognition step concentrates this risk. A copilot that pipelines raw
LLM output into a planner and executor must, at minimum, emit a well-formed,
schema-conformant, authorised intent structure before any downstream component
acts on it. Unless that structured output is validated deterministically, the
entire framework inherits the model's failure modes. This dissertation addresses
the resulting problem: \emph{how can a small, offline, CPU-only language model be
made sufficiently reliable to serve as the intent-recognition head of an
automated security pipeline, when its raw output is constrained by deterministic
validation rather than by a human reviewer or a cloud back-end?}